\subsection{Algebra: geometric regularity and Cohen--Macaulay field extensions}\label{algebra-proof-071-algebra-geometric-regularity-and-cohenmacaulay-field-extensions}

This editorial verification covers the original Stacks Project \texttt{algebra.tex} 47092--47318 at authority \texttt{a044\allowbreak{}46e5\allowbreak{}7ec1\allowbreak{}fbc2\allowbreak{}52a8\allowbreak{}71af\allowbreak{}cec7\allowbreak{}752f\allowbreak{}b280\allowbreak{}7b14}, SHA-256 \texttt{FA8B\allowbreak{}B92E\allowbreak{}58A4\allowbreak{}F78A\allowbreak{}2BD0\allowbreak{}1B3B\allowbreak{}6A4A\allowbreak{}87DE\allowbreak{}0A0D\allowbreak{}279F\allowbreak{}5DD9\allowbreak{}0641\allowbreak{}B574\allowbreak{}DD5F\allowbreak{}BFFF\allowbreak{}A4F3}. The \href{https://github.com/stacks/stacks-project/blob/a04446e57ec1fbc252a871afcec7752fb2807b14/algebra.tex\#L47092}{original source} remains identifiable. The two reports retain MC-\AlgebraIdentifierBreak{}STK-\AlgebraIdentifierBreak{}ERR-\AlgebraIdentifierBreak{}0738, changing only the left-hand field in 47214. No additional source operation is proposed. All arguments here are editorial.

Original source dependencies read in this review include 6165--6181, 27488--27503, 39289--39335, 46386--46400, 46609--46626 and 46714--46732. The completed field-formal-smoothness, geometric-normality, ascent, descent and dimension-map notes supply their already proved comparisons. Three new indexed routing hits were returned; none of their contents was read. This is not a new literature-reading or novelty claim.

\subsubsection{Noetherianity and the original regularity tests}\label{algebra-proof-071-noetherianity-and-the-original-regularity-tests}

The source uses a Noetherian regular ring, with every prime localization a regular local ring. It tests geometric regularity against finitely generated field extensions, or equivalently finite purely inseparable ones. These quantifiers remain unchanged. The zero ring is Noetherian and has no prime localizations; all scalar extensions of it are zero and satisfy the same conventions.

Let \(R\) be the original Noetherian \(k\)-algebra and let \(K/k\) be finitely generated as a field. Choose field generators \(u_1,\ldots,u_n\), let \(D=k[u_1,\ldots,u_n]\subset K\), and put \(U=D\setminus\{0\}\). The map \[
 (U^{-1}D)\otimes_kR\longrightarrow(U\otimes1)^{-1}(D\otimes_kR),
 \quad(d/u)\otimes r\longmapsto(d\otimes r)/(u\otimes1)
\] has inverse \((\sum_jd_j\otimes r_j)/(u\otimes1)\mapsto\sum_j(d_j/u)\otimes r_j\). The fraction equality witnesses and tensor relations establish both composites. The ring \(D\otimes_kR\) is a quotient of the original finite polynomial ring \(R[X_1,\ldots,X_n]\), hence Noetherian. A localized ideal is generated by the images of a finite generating set of its contraction, so its localization is Noetherian. Since \(U^{-1}D=K\), this proves every Noetherianity assertion needed for the finitely generated scalar extensions.

For the two regularity tests, one implication is inclusion of the quantified classes. Conversely let the finite purely inseparable test hold, and fix the source's finitely generated \(K/k\). The already proved make-separable diagram has \(k'/k\) and \(K'/K\) finite purely inseparable, with \(K'/k'\) separable. The geometric-normality note retains its actual multiplication map \(k'\otimes_kK\to k'K\), nilpotent kernel and original field embeddings. We do not replace this possibly nonreduced tensor by a field without that quotient map.

Let \(B\) be the original smooth \(k'\)-domain with fraction field \(K'\), supplied by the completed generic smooth-model proof. The ring \(k'\otimes_kA\) is regular by the assumed finite test. Smooth base change gives \[
 k'\otimes_kA\longrightarrow B\otimes_{k'}(k'\otimes_kA),
\] so regular smooth ascent, proved in the ascent note, makes the target regular. The exact comparison is \[
 b\otimes(\lambda\otimes a)\longmapsto b\lambda\otimes a,
 \qquad b\otimes a\longmapsto b\otimes(1\otimes a).
\] It identifies the target with \(B\otimes_kA\). Localizing at every original nonzero \(b\in B\) gives \(K'\otimes_kA\); the inverse fraction maps are those displayed above. Each prime localization is an original regular local ring, so this localization remains regular and Noetherian. Finally \[
 K\otimes_kA\longrightarrow K'\otimes_kA
\] is faithfully flat: for an original \(K\otimes_kA\)-module \(M\), its base change identifies with \(K'\otimes_KM\) by \((u\otimes a)\otimes m\mapsto u\otimes am\). The inverse sends \(u\otimes m\) to \((u\otimes1)\otimes m\). A nonempty field basis proves exactness and reflection of zero. The complete regular-descent argument in the descent note now proves the original \(K\otimes_kA\) regular. This proves the source equivalence without an unrestricted colimit assertion about regular rings.

If \(A\) is geometrically regular and \(K/k\) is finitely generated, then \(K\otimes_kA\) is geometrically regular over \(K\): a finitely generated \(L/K\) is finitely generated over \(k\), and \[
 L\otimes_K(K\otimes_kA)\cong L\otimes_kA,
 \quad \ell\otimes(u\otimes a)\longleftrightarrow\ell u\otimes a
\] has inverse \(\ell\otimes a\mapsto\ell\otimes(1\otimes a)\). Noetherianity of the base-changed algebra was proved above. The later exponent-one criterion mentioned at 47160--47163 is only a forward reference to More on Algebra, not a new theorem asserted or re-proved in this intake.

For the source's faithfully flat map \(A\to B\) with geometrically regular \(B\), first descend Noetherianity. If \(I\subset A\), finitely many original elements of \(I\) generate \(IB\); let their ideal be \(I_0\). Flatness identifies \((I/I_0)\otimes_AB\) with \(IB/I_0B=0\). Faithfulness gives \(I=I_0\). Thus every ideal of \(A\) is finite. After any finite purely inseparable \(L/k\), the map \(A\otimes_kL\to B\otimes_kL\) is faithfully flat by the same base-change module comparison, and the target is regular. Regular descent proves the source test. For a smooth \(A\to B\) with geometrically regular \(A\), finite presentation makes \(B\) Noetherian. Every finitely generated scalar extension of that smooth map is smooth, and regular smooth ascent proves the source assertion. These arguments also cover the zero algebras whenever allowed by the original maps.

\subsubsection{Finite purely inseparable descent to an actual subfield}\label{algebra-proof-071-finite-purely-inseparable-descent-to-an-actual-subfield}

Retain the original directed union \(k=\bigcup_i k_i\) and the finite purely inseparable extension \(k'/k\). The construction below proves the field and tensor identity used at 47213--47215. It does not assume that an arbitrary descended finite algebra is a field.

In characteristic \(p>0\), choose a \(k\)-basis \(e_1=1,e_2,\ldots,e_d\) of the original \(k'\). There is an exponent \(q=p^N\) such that every \(e_j^q\) belongs to \(k\). Retain all multiplication constants \[
 e_re_s=\sum_{j=1}^d c_{rsj}e_j,\qquad c_{rsj}\in k,
\] and all elements \(a_j=e_j^q\in k\). Directedness puts these finitely many original constants and elements in one \(k_i\). Define the actual subspace of \(k'\) \[
 k_i'=\left\{\sum_{j=1}^d b_je_j:b_j\in k_i\right\}.
\] It contains \(1\), is closed under addition and multiplication by the retained constants, and is a domain as a subring of the original field. Its displayed basis is linearly independent over \(k_i\), since it is independent over \(k\). For any nonzero \(x\in k_i'\), multiplication by \(x\) is an injective endomorphism of this finite-dimensional \(k_i\)-space, hence surjective. In particular \(xy=1\) for some \(y\in k_i'\). Thus \(k_i'\) is a field of degree \(d\) over \(k_i\).

For every original element \(x=\sum_jb_je_j\) of this field, Frobenius gives the full formula \[
 x^q=\sum_{j=1}^d b_j^q e_j^q=\sum_{j=1}^d b_j^q a_j\in k_i.
\] The mixed binomial coefficients vanish in characteristic \(p\). This proves that \(k_i'/k_i\) is finite purely inseparable, with the actual bound \(q\); none of its coefficients has been discarded.

The multiplication comparison \[
 \mu:k\otimes_{k_i}k_i'\xrightarrow{\sim}k',\qquad
 \lambda\otimes x\longmapsto\lambda x
\] sends the \(k\)-basis \(1\otimes e_j\) to the original basis \(e_j\). It is therefore bijective and multiplicative, and its inverse is explicitly \[
 \sum_j\lambda_je_j\longmapsto\sum_j\lambda_j\otimes e_j.
\] The retained multiplication constants verify multiplication in this inverse formula as well. In characteristic zero, a finite purely inseparable extension is \(k'=k\); choose any index, set \(k_i'=k_i\), and use the identity multiplication comparison. The degree-one case in positive characteristic is included too.

Since \(A\) is an algebra over the original \(k\), the needed algebra comparison is \[
 \Phi:A\otimes_{k_i}k_i'\xrightarrow{\sim}A\otimes_kk',
 \qquad a\otimes x\longmapsto a\otimes x.
\] For \(y=\sum_j\lambda_je_j\in k'\), its inverse sends \[
 a\otimes y\longmapsto\sum_j a\lambda_j\otimes e_j.
\] The unique basis expansion proves additivity; the \(k\)-balanced relation moves any coefficient from the second tensor factor to the first in precisely this formula. The multiplication constants prove multiplicativity and both composites fix each elementary tensor. This is exactly the source's scalar-extension identity with the corrected left-hand field \(k'\).

Noetherianity of the unchanged \(A\) follows from its geometric regularity over any one of the nonempty directed family of fields. Its geometric regularity over the selected \(k_i\) makes \(A\otimes_{k_i}k_i'\) regular. Transporting along the proved \(\Phi\) verifies every finite purely inseparable test over \(k\), and hence geometric regularity over \(k\).

The original printed \(k_i=k\otimes_{k_i}k_i'\) is not this statement. For example, take the constant directed system \(k_i=k=\mathbf F_p(t)\) and \(k'=k(t^{1/p})\). The right side of the corrected comparison has degree \(p\) over \(k\), whereas \(k_i=k\) has degree one; irreducibility follows from Eisenstein at \(t\). The two reports therefore identify one mathematical correction, already present as unit 0738, not two operations.

\subsubsection{Separable algebraic change of the ground field}\label{algebra-proof-071-separable-algebraic-change-of-the-ground-field}

Let \(k'/k\) be the original separable algebraic extension and \(A\) a \(k'\)-algebra. Noetherianity is a property of the same ring \(A\) in either statement. If \(L/k\) is finite purely inseparable, then \(L'=k'\otimes_kL\) is a field by the reducedness and Frobenius-unit proof retained in the geometric-normality note. It is finite of degree \([L:k]\) over \(k'\) and purely inseparable: for a common exponent \(q\), every \(\sum_jb_j\otimes\ell_j\) has \(q\)-th power \(\sum_jb_j^q\ell_j^q\) in the embedded \(k'\). The exact associativity isomorphism \[
 A\otimes_{k'}(k'\otimes_kL)\longrightarrow A\otimes_kL,
 \quad a\otimes(b\otimes\ell)\longmapsto ab\otimes\ell
\] has inverse \(a\otimes\ell\mapsto a\otimes(1\otimes\ell)\). Hence geometric regularity over \(k'\) implies the original finite test over \(k\).

For the converse, first suppose \(k'/k\) finite separable. The preceding finitely generated base-change argument makes \(A\otimes_kk'\) geometrically regular over its second \(k'\)-factor. The multiplication map \(A\otimes_kk'\to A\) is an étale localization. Here are its exact denominators. Choose a primitive element \(\alpha\) of \(k'/k\) and its original monic minimal polynomial \(P(T)\). In \(k'[T]\), write \[
 P(T)=(T-\alpha)Q(T),\qquad Q(\alpha)=P'(\alpha)\ne0.
\] The map \(A[T]/(P)\to A\otimes_kk'\), \(\sum_j a_jT^j\mapsto\sum_j a_j\otimes\alpha^j\), is an isomorphism by the original power basis. On the left, \(\alpha\) in a coefficient means its given image in \(A\). Inverting \(Q\) forces \(T=\alpha\), so evaluation gives \[
 (A[T]/(P))_Q\xrightarrow{\sim}A,
 \quad f(T)/Q(T)^n\longmapsto f(\alpha)/Q(\alpha)^n,
\] with inverse the coefficient inclusion. The element \(Q(\alpha)\) is already a unit in \(A\), including when \(A=0\), so all denominators are defined. This proves that the multiplication map is localization at the single indicated element, hence is étale by 39293 and smooth. Evaluation sends \(1\otimes r\) to the original \(r\in A\), so it preserves exactly the second \(k'\)-algebra structure. Geometric regular smooth ascent proves the desired assertion over \(k'\).

For arbitrary separable algebraic \(k'/k\), its finite intermediate extensions form a directed union. The argument just given proves that the same \(A\) is geometrically regular over each of them. The finite purely inseparable descent construction in the preceding section applies to this union and gives geometric regularity over the original \(k'\). It is the finite-test descent, not a claim that every infinite scalar extension of a Noetherian ring stays Noetherian.

\subsubsection{The original field tensor and its Cohen--Macaulay localizations}\label{algebra-proof-071-the-original-field-tensor-and-its-cohenmacaulay-localizations}

Let \(K/k,L/k\) be the original field extensions, one finitely generated. The tensor flip \(x\otimes y\mapsto y\otimes x\) is its own inverse, so suppose the source's \(K/k\) is finitely generated. Choose its original transcendence basis \(t_1,\ldots,t_r\); then \(K\) is finite of degree \(d\ge1\) over \(E=k(t_1,\ldots,t_r)\). Set \[
 U=k[t_1,\ldots,t_r]\setminus\{0\},\qquad
 R=E\otimes_kL\cong U^{-1}L[t_1,\ldots,t_r],
 \quad S=K\otimes_kL.
\] The displayed localization inverts the original nonzero polynomials with coefficients in \(k\), not every nonzero polynomial with coefficients in \(L\). The polynomial algebra over the field \(L\) is Cohen--Macaulay by the preceding source polynomial theorem, and so is each of its localizations. Thus \(R\) is Noetherian and Cohen--Macaulay, including the case \(r=0\), when \(R=L\).

The map \[
 K\otimes_E(E\otimes_kL)\xrightarrow{\sim}K\otimes_kL,
 \quad x\otimes(e\otimes\ell)\longmapsto xe\otimes\ell
\] has inverse \(x\otimes\ell\mapsto x\otimes(1\otimes\ell)\). An original \(E\)-basis of \(K\) therefore makes \(S\) finite free of rank \(d\) over \(R\). It is Noetherian, and this finite map is integral. For completeness, multiplying an element of \(S\) on its finite free \(R\)-module gives a matrix; its monic characteristic polynomial annihilates that matrix, and evaluation on the unit gives the integral equation for the original element.

Take any original prime \(\mathfrak q\subset S\), with contraction \(\mathfrak p\subset R\). Then \(R_{\mathfrak p}\to S_{\mathfrak q}\) is a flat local map of Noetherian local rings. We do not assume that localization at this one prime preserves finiteness as a module. Instead, any strict chain of primes in \(S\) below \(\mathfrak q\) contracts to a strict chain below \(\mathfrak p\), by incomparability for the original integral extension. Thus \[
 \dim S_{\mathfrak q}\le\dim R_{\mathfrak p}.
\] Let \(n=\dim R_{\mathfrak p}\). A parameter regular sequence of length \(n\) in the Cohen--Macaulay \(R_{\mathfrak p}\) remains regular on \(S_{\mathfrak q}\): flatness preserves every multiplication injection and quotient, and the local map is faithfully flat, so the final nonzero quotient remains nonzero. Consequently \[
 n\le\operatorname{depth}S_{\mathfrak q}
   \le\dim S_{\mathfrak q}\le n.
\] All inequalities are equalities. The target local ring is Cohen--Macaulay; since the prime was arbitrary, this proves the original tensor theorem. It uses the second alternative of 27492--27493 when needed, without an unjustified local finiteness assertion or any separability hypothesis on \(K/E\).

\subsubsection{The exact residue-fibre comparison and the earlier consequence}\label{algebra-proof-071-the-exact-residue-fibre-comparison-and-the-earlier-consequence}

Retain the final source notation \(S_K=K\otimes_kS\), \(\mathfrak q\subset S\) and \(\mathfrak q_K\subset S_K\) lying over it. Put \(A=S_{\mathfrak q}\), \(B=(S_K)_{\mathfrak q_K}\), and \(F=K\otimes_k\kappa(\mathfrak q)\). These are the original rings; \(A\to B\) is flat local, and \(A,B\) are Noetherian by the first section. Write \[
 D=K\otimes_kS_{\mathfrak q},\quad
 P=\mathfrak q_KD,\quad
 J=K\otimes_k\mathfrak q S_{\mathfrak q}.
\] The ideal \(P\) exists because every denominator from \(S\setminus\mathfrak q\) avoids \(\mathfrak q_K\); it is prime and contains \(J\). The exact quotient \(D/J\cong F\) sends \(u\otimes(a/v)\) to \(u\otimes\overline{a/v}\). Let \(\mathfrak q'\) be the image of \(P\) in \(F\). Quotient and localization then give \[
 B/\mathfrak q B\cong(D/J)_{P/J}=F_{\mathfrak q'}.
\] The forward map takes the residue class of \(z/w\in D_P\) to \(\bar z/\bar w\). Its inverse chooses lifts of numerator and denominator in \(D\); two choices differ by elements of \(J\), and the same fraction equality witnesses prove independence after quotienting. A denominator outside \(P/J\) has every lift outside \(P\), so both formulas are defined. Composing with the tensor flip identifies this with exactly the source's \((\kappa(\mathfrak q)\otimes_kK)_{\mathfrak q'}\), relabeling the prime by that proved isomorphism.

The field tensor just proved is Cohen--Macaulay; thus its localization \(B/\mathfrak q B\) is Cohen--Macaulay. The earlier complete flat local formulas, retained in the ascent note, are \[
 \dim B=\dim A+\dim(B/\mathfrak qB),\qquad
 \operatorname{depth}B=\operatorname{depth}A+
                      \operatorname{depth}(B/\mathfrak qB).
\] These finite local dimensions are not global bounds on the original rings. Subtracting the two full formulas gives \[
 \dim B-\operatorname{depth}B=
 \dim A-\operatorname{depth}A+
 \dim(B/\mathfrak qB)-\operatorname{depth}(B/\mathfrak qB).
\] The final summand is zero by the proved fibre theorem. Hence \(A\) is Cohen--Macaulay exactly when \(B\) is, proving the original final lemma at each specified pair of primes.

This review also resolves the scope of earlier group 650. Its proof at 27488--27503 extended the particular zero-dimensional-fibre assertion to a flat local map with any Cohen--Macaulay closed fibre. But the \href{https://github.com/stacks/stacks-project/blob/a04446e57ec1fbc252a871afcec7752fb2807b14/algebra.tex\#L46386}{later original lemma at 46386--46400} explicitly states the full equivalence, and the earlier dimension formula supplies the same dimension equality. Thus that broader implication is already in the original work. Preserve group 650 and its full regular-sequence proof for provenance, while refining its current ledger classification to a consequence already stated in later source content. It is a valid receiving proof, not an extension beyond the entire source. This comparison propagates to the final field-extension application just proved.

\subsubsection{Report accounting and continuation}\label{algebra-proof-071-report-accounting-and-continuation}

OCC-01032 and OCC-12449 both retain the single existing operation MC-\AlgebraIdentifierBreak{}STK-\AlgebraIdentifierBreak{}ERR-\AlgebraIdentifierBreak{}0738-\AlgebraIdentifierBreak{}OP1 at 47214. The basis construction proves the corrected field and both scalar-extension maps with all original constants. No new source operation, conceptual-extension group, external reading, build or publication is claimed. Full proofs and the refinement of group 650 stay editorial. Continue source review at 47319; the remaining thirteen reports have been read, but their proof intervals have not yet been adjudicated.
